\label{inv:C3.2}\label{wb:lem:bcpr-degree-one-input}
\StatementRecord{inv:C3.2}{proved}{}{the hypotheses, notation and definitions stated here; \cref{inv:eq:C-collar-sum}; \cref{inv:C3.1}}

Use the marked cohomology class $\gamma$ dual to
$\hat\gamma$, the analytic peripheral generators and the lift relation of
[AC26, Section 7.5], and put
$p=12l_0-4l_1-3l_2=\pm1$.  Let $g_\gamma$ be the unique coherent
Leray class whose fibre image is $\gamma^{0,1}$.  Then
\begin{equation}\label{inv:eq:C-factor12}
\exp(12g_\gamma)=f^*\OO_{\Pone}(-p).
\end{equation}
For the standard data $(l_0,l_1,l_2)=(0,1,-1)$, this is
$\exp(12g_\gamma)=f^*\OO(1)$.  For the opposite value $p=1$,
the exponent representing $f^*\OO(1)$ is $-12g_\gamma$.
\end{proposition}


\begin{proof}
\begin{enumerate}[label=\arabic*.,leftmargin=*]
\item We include the comparison that fixes the sign.  A degree-one integral
class $k\gamma$ on a regular fibre extends over the cyclic filling of
order $m_j$ precisely when the affine generator $\hat g_j$ can be
assigned an integer value $n_j$ satisfying
$m_j n_j=k\gamma(v_j)=kl_j$.  Admissibility gives
$\gcd(l_j,m_j)=1$, so this is precisely $m_j\mid k$.
At the cusp $\gamma$ extends, as follows from its actual specialization
to $H^1(W;\Z)$ in [AC26, Theorem B.1 and Proposition 7.12].
The global monodromy invariants have rank one generated by $\gamma$.
It follows, by intersection of these three local images in the marked
local system, that
$H^0(B,R^1f_*\Z)=12\Z\gamma$.  This recovers the degree-one
part of [AC26, Proposition 7.26], without identifying the reduced and
non-reduced fibre: this is topological cohomology, unchanged by
nilpotents, whereas the coherent comparison uses the full fibre.

\item Let $U=\bigsqcup_iD_i$ be small discs around the special points and
let $V$ be their pair-of-pants complement with collars.  Lift
$12\gamma$ to integral classes $\Theta_i$ on $f^{-1}(D_i)$ and
$\Theta_V$ on $f^{-1}(V)$.  Existence over $V$ follows from
$H^2(V;\Z)=0$ and the low-degree Leray sequence; over each disc it
follows from the local images just computed.  Write
$a_j=\Theta_V(\hat\rho_j)$.  On the two finite analytic peripheral generators
the differences $\Theta_i-\Theta_V$ have periods
\[
n_1=4l_1-a_1,\qquad n_2=3l_2-a_2.
\]
At the cusp, the specified lift relation
$\widetilde a_0=t_\mu\hat\rho_2^{-1}\hat\rho_1^{-1}$,
with $\gamma(\mu)=l_0$, and the null-homotopy of
$\widetilde a_0$ in the cusp filling give
$n_0=-12l_0+a_1+a_2$.  Hence
\begin{equation}\label{inv:eq:C-collar-sum}
n_0+n_1+n_2=-p.
\end{equation}
These are the explicitly oriented boundary evaluations in [AC26,
Proposition 7.27, pp. 62--63], not an unspecified sign of a spectral
sequence differential.

\item Represent the lifts by real closed one-forms.  On a collar their
difference is the pullback of a real closed base form $b_i$ with
clockwise period $n_i$, modulo an exact form.  This follows from the
degree-one Leray sequence over the annulus, since the difference has
zero fibre restriction and the annulus has no second cohomology.
Subtract exact terms on smaller collars, extending their primitives
with cutoffs, so the equality of forms holds wherever a partition of
unity changes.  For a real base cutoff $\rho_i$ equal to one toward
$D_i$ and zero toward $V$, the resulting glued real one-form $\alpha$
on $X$ has
\[
d\alpha=f^*\vartheta,\qquad
\vartheta=\sum_i d\rho_i\wedge b_i.
\]
The form $\vartheta$ is a smooth real base two-form; thus
$\bar\partial\alpha^{0,1}=0$.  Its fibre restriction is the
$(0,1)$ image of $12\gamma$, so the coherent Leray edge isomorphism
$H^1(X,\OO_X)\simeq H^0(B,R^1f_*\OO_X)$ identifies
$[\alpha^{0,1}]=12g_\gamma$.

\item In a positively oriented polar coordinate $(r,\varphi)$ on $D_i$,
$b_i=-n_i\,d\varphi/(2\pi)$ modulo an exact form.
Because $\rho_i$ decreases from one to zero as $r$ increases,
\[
\int_{D_i\cap V}d\rho_i\wedge b_i
=-n_i\int_0^1d\rho_i=n_i.
\]
The exact part does not change this integral, since $d\rho_i$ is
supported inside the collar.  Therefore
$\int_B\vartheta=-p$ by \eqref{inv:eq:C-collar-sum}.  \ref{inv:C3.1} identifies
the Bott--Chern class of $\exp(12g_\gamma)$ with $[f^*\vartheta]$.
On $\Pone$ two real $(1,1)$ forms of the same integral differ by
$i\partial\bar\partial$ of a real function (solve the ordinary
scalar Laplace equation).  Hence this class is the Chern class of
$f^*\OO(-p)$, and injectivity in \ref{inv:C3.1} proves
\eqref{inv:eq:C-factor12}.  The program
\texttt{code/c\_sign\_and\_fibre.py} verifies the integer
intersection, \eqref{inv:eq:C-collar-sum}, and the collar integral exactly.
\end{enumerate}
\end{proof}